% Standalone mathematical section for the root report. No preamble required.
\section{Five inference patterns with closure clauses}
\label{sec:clause-five}

This construction combines three already established ingredients: the
canonical successor-seeded recursor, one-term crossed transport, and the
older clause presentation of the paired binder. The combination removes
the first-class $Q$ constructor and merges its specialization rule with
ordinary assumption use. Its novelty is the resulting complete interface
and comparison; none of those three ingredients is claimed as new.

\subsection{Syntax and every primitive rule}

Numerical terms are canonical finite expressions
\[
 t::=x\mid0\mid R(n,a;u.h),\qquad
 R(0,a;u.h)\equiv R(0,a;u.0).
\]
Only the displayed $u$ in $h$ is bound. Put
$Sa:=R(0,a;u.0)$. Every substitution is capture-avoiding substitution
followed by the canonicalization that erases the body of a literal
zero-counter recursor. Thus all occurrence tests below concern canonical
syntax. The formula grammar is only
\[
 F::=(s=t)\mid D_{x,y}F.
\]
A reusable context is a finite list of clauses
\[
 [A;y.B(y)],\qquad y\notin\operatorname{FV}(A).
\]
Clauses are context objects, not formulas or public goals. Their ordinary
reading is $A\to\forall yB(y)$. Let
\[
 J(a):=R(S0,0;u.a),\qquad c:=J(0),\qquad H:=(c=0).
\]
The update variable in $J(a)$ is absent from $a$. A retained ordinary
assumption $F$ means the constant clause $[H;v.F]$ with $v$ unused.
This abbreviation applies also to the local assumption in the following
introduction rule.

There are exactly the following five primitive inference patterns.
\begin{enumerate}
\item Computation, in any context:
\[
 R(Sn,a;u.h)=h\langle R(n,a;u.h)/u\rangle.
 \tag{R}
\]
\item One-term crossed transport, with unrestricted source:
\[
 \frac{\Gamma\vdash s=t\qquad\Gamma\vdash h(s)=r}
      {\Gamma\vdash r=h(t)}.
 \tag{T}
\]
Here $h$ is a single fixed term template; constant templates are allowed.
\item Clause use:
\[
 \frac{\Gamma\vdash A}{\Gamma\vdash B(t)}
 \quad\text{if }[A;y.B(y)]\in\Gamma.
 \tag{U}
\]
\item Paired elimination:
\[
 \frac{\Gamma\vdash D_{x,y}B(x,y)\qquad
       \Gamma,[B(b,0);y.B(b,y)]\vdash F}
      {\Gamma\vdash F}.
 \tag{E}
\]
The parameter $b$ is fresh for the retained context, conclusion, and
fixed parameters of $B$.
\item Witness-directed introduction:
\[
 \frac{\Gamma,B(t,\eta)\vdash B(t,Sz)}
      {\Gamma\vdash D_{x,y}B(x,y)},\qquad
 \eta=\begin{cases}
 z,&B\text{ is a literal equation and }t\equiv0,\\
 0,&\text{otherwise}.
 \end{cases}
 \tag{I}
\]
The parameter $z$ is fresh for the retained context, witness $t$, and
fixed parameters of $B$.
\end{enumerate}
All displayed instances use normalized substitution. There is no separate
hypothesis rule, zero computation rule, $Q$ formula constructor, $Q$
elimination rule, reflexivity axiom, numerical case rule, or induction
rule. These statements describe the interface, not an uncharged saving:
the clause pair, its formal parameter, and its substitution are part of
the implementation.

\subsection{The elementary derivations}

Computation proves $J(a)=a$ uniformly in every context, and in particular
proves $H$. Ordinary assumption use is therefore computation of $H$
followed by U on the constant clause. Weakening, fresh renaming, and
ordinary-assumption cut follow by induction on derivations. In a cut,
each use of the distinguished constant clause is replaced by the supplied
proof of its constant conclusion. No general numerical-substitution
theorem for the syntax-directed I rule is presupposed here.

\begin{lemma}[Equality from one term template]
The five-rule calculus derives reflexivity, symmetry, transitivity,
numerical congruence, and transport in every literal-equation template.
\end{lemma}
\begin{proof}
Use $J(a)=a$ as both premises of T with the identity template to get
$a=a$. To reverse a supplied $s=t$, use $H$ as major premise and the
constant template $h(v)=s$; T returns $t=s$. Given $s=t$ and $t=r$,
reverse the former and apply T with identity template to major $t=r$
and minor $t=s$, obtaining $s=r$. Congruence uses major $s=t$ and
reflexivity of $h(s)$.

For $E(v)\equiv a(v)=b(v)$, the two applications
\[
 s=t,\ a(s)=b(s)\ \Longrightarrow\ b(s)=a(t),
 \qquad
 s=t,\ b(s)=a(t)\ \Longrightarrow\ a(t)=b(t)
\]
give arbitrary equational transport. Each term template is fixed finite
syntax; no function variable or function extensionality is introduced.
\end{proof}

\begin{lemma}[All case introductions]
For every formula $B$ and every term $t$, the case rule
\[
 \frac{\Gamma,B(t,0)\vdash B(t,Sz)}
      {\Gamma\vdash D_{x,y}B(x,y)}
\]
is admissible.
\end{lemma}
\begin{proof}
The only nonprimitive instance has an equation matrix $E$ and literal
zero witness. Use the nonliteral term $c=J(0)$ instead. Under $E(c,0)$,
the theorem $c=0$ and equational transport give $E(0,0)$. Insert the
supplied branch by ordinary-assumption cut, obtaining $E(0,Sz)$, and
transport back to $E(c,Sz)$. Primitive I at witness $c$ is the case mode
and proves the original goal $D_{x,y}E(x,y)$. The goal has not changed;
only the witnessing term used by its proof has changed.
\end{proof}

\begin{lemma}[Transport in every formula]
For every Eq/$D$ template $F(v)$, from $s=t$ and $F(s)$ one can derive
$F(t)$, under arbitrary retained clauses.
\end{lemma}
\begin{proof}
Induct externally on the finite formula template. The equation case was
proved above. For $F(v)=D_{x,y}B(v,x,y)$, eliminate $F(s)$ at fresh $b$,
obtaining the clause $[B(s,b,0);y.B(s,b,y)]$. Introduce the target by the
derived case rule with witness $b$. Convert its local $B(t,b,0)$ to
$B(s,b,0)$ using symmetry and the induction hypothesis, apply U at $Sz$,
then convert $B(s,b,Sz)$ to $B(t,b,Sz)$ using the induction hypothesis.
Close introduction and elimination. Only the strictly smaller matrix is
used in the recursive transport, so the construction is not circular.
\end{proof}

\subsection{Closure as an abbreviation, not a constructor}

For a fresh unused variable $w$, write
\[
 q_yB(y):=D_{w,y}B(y).
\]
The formula has the following three derived interfaces.
\begin{itemize}
\item From $\Gamma,B(0)\vdash B(Sz)$ derive $\Gamma\vdash q_yB$ by
the derived case introduction with unused witness.
\item From $q_yB$ and $B(0)$ derive $B(t)$: eliminate the paired formula,
use the opened clause at $t$, and close elimination. Its witness is unused.
\item In any context containing $[B(0);y.B(y)]$, derive $q_yB$: under
its case-introduction assumption $B(0)$ use that clause at $Sz$.
\end{itemize}
These are respectively introduction, elimination, and eta expansion of a
closure resource. The third is what permits translation of a native DE
branch that formerly assumed a first-class $Q$ formula.

For an equation $E(y)$, primitive I at witness zero and unused witness
coordinate gives $q_yE$ directly from
$\Gamma,E(z)\vdash E(Sz)$. Combining this with the derived closure
elimination gives ordinary contextual induction on literal equations.
Combining closure introduction with a separately supplied base gives
numerical cases for arbitrary native formulas. All progress introductions
at a nonzero or nonliteral original witness can also be recovered by the
usual unused-coordinate construction: form $q_yE(t,y)$ by progress, then
under the base $E(t,0)$ specialize its closure at $Sz$ and use case
introduction at $t$.

\subsection{Exact comparison with the six-rule canonical calculus}

Let $\mathcal K_6$ be the repository's canonical successor-seeded Eq/$Q$/$D$
calculus, including its restricted crossed equational transport and
root-reversing QE. The canonical-recursors equivalence theorem already
gives it ordinary equality transport and ordinary QE as derived rules,
and both forms of paired introduction. The following comparison changes
neither its term syntax nor its normalization.

Define $T:\mathcal K_6\to\mathcal K_5$ by
\[
 T(s=t)=(s=t),\qquad
 T(D_{x,y}B)=D_{x,y}T(B),\qquad
 T(Q_yB)=q_yT(B).
\]
Every newly inserted dummy coordinate is fresh and unused. This finite
translation preserves free numerical variables, the equation/compound
root distinction, and normalized numerical substitution.

\begin{theorem}[Preservation in both directions]
For every native context $\Gamma$ and formula $A$,
\[
 \Gamma\vdash_{\mathcal K_6} A
 \quad\Longrightarrow\quad
 \operatorname{ord}(T\Gamma)\vdash_{\mathcal K_5}T(A).
\]
Conversely, for Eq/$D$ formulas $\Delta,F$,
\[
 \operatorname{ord}(\Delta)\vdash_{\mathcal K_5}F
 \quad\Longrightarrow\quad
 \Delta\vdash_{\mathcal K_6}F.
\]
Here $\operatorname{ord}(\Delta)$ replaces each ordinary assumption by its
constant clause $[H;v.F]$.
\end{theorem}
\begin{proof}
Forwards, induct on the $\mathcal K_6$ derivation. Hyp uses $H$ and the
constant clause. Computation is unchanged. Its equation transport and
root reversal are supplied by derived equational transport and symmetry.
QE becomes the derived $q$ elimination, followed by symmetry when its
conclusion is an equation. The unified introduction is preserved exactly:
the translation preserves roots and its numerical witness is unchanged.

For DE, the translated premise proves $D_{x,y}T(B)$. Its elimination
opens $[T(B)(b,0);y.T(B)(b,y)]$. Eta expansion derives the translated old
branch assumption $q_yT(B)(b,y)$. Ordinary-assumption cut inserts this
proof into the inductively translated branch. Close DE. All new dummy
variables can be freshened before the construction.

Backwards, track each context occurrence by the rule that introduced it.
An ordinary constant clause is represented by its ordinary formula
assumption; a clause opened by E is represented by the corresponding
$Q_yB(b,y)$ assumption. U on an ordinary clause is Hyp; U on an opened
closure is ordinary QE. Computation, T, E, and I are respectively the
same computation, derived one-term equality transport, native DE, and
the same witness-directed introduction. This proves preservation by
induction, including derivations from arbitrary ordinary contexts.

If one also wants a statement for arbitrary externally supplied clauses
$[A;y.B]$, interpret each by the already derived ordinary macro
$A\to\forall yB(y)$. U is then implication elimination followed by
universal elimination. No externally supplied clause is needed by a
closed-certificate verifier.
\end{proof}

\begin{theorem}[Both composites]
The inclusion $U$ of Eq/$D$ formulas into the native Eq/$Q$/$D$ language
satisfies $TU(F)=F$ literally. For every native formula $A$,
\[
 A\dashv\vdash_{\mathcal K_6}UT(A).
\]
Both numerical translations are identities.
\end{theorem}
\begin{proof}
The literal composite follows immediately from the translation's clauses.
For the other composite use the native equivalence
\[
 Q_yB(y)\dashv\vdash D_{w,y}B(y)\quad(w\text{ unused}).
\]
Forwards, under $B(0)$ specialize the supplied $Q$ at $Sz$ and introduce
the paired formula by cases. Backwards, DE opens exactly the desired
$Q_yB$ at an unused witness, and Hyp closes its branch. Formula congruence
then lifts this equivalence recursively through $Q$ and $D$. Thus
external structural induction proves the displayed native composite.
No quotient, retract, standard-number restriction, or reinterpretation
of numerical equality is introduced.
\end{proof}

Composing this explicit equivalence with the repository's existing
canonical-recursor bi-interpretation yields equiconsistency and
bi-interpretability with full first-order iPRA. Both numerical domains
are whole, equality is actual equality, and both numerical composite
graphs remain $x=y$. The designated contradiction $0=S0$ is unchanged.

\subsection{What the smaller interface does and does not establish}

The public formula constructors decrease from equality/$Q$/$D$ to
equality/$D$ and the named inference patterns decrease from six to five.
The implementation must still store a closure antecedent and a one-bound-
variable body in each context entry. Ordinary assumptions require a
neutral computation premise, so certificates can become longer even
when the verifier becomes shorter. A fair size comparison must count
canonical substitution, clause formation and opening, scope validation,
input parsing, and the public-goal validator in both programs.

The mathematical equivalence above is a written proof. It is not itself
a new Lean formalization, a proof of a global lower bound, or evidence
that every historical compact checker has been beaten. Those are
separate claims requiring their own measured or formal evidence.
